\section*{Problem 1}

设矩阵
\[
A = \begin{bmatrix}
    1 & 2 \\
    2 & 4 \\
\end{bmatrix}.
\]
\begin{enumerate}
    \item 求 $A$ 的秩，以及列空间 $C \left(A\right)$ 与零空间 $N \left(A\right)$ 各自的一组基。
    \item 对 $A$ 进行 SVD，在该矩阵的 SVD 中，说明非零奇异值对应的左、右奇异向量分别张成哪个基本子空间；补齐后的其余左、右奇异向量又分别张成哪个基本子空间。
    \item 将向量 $\boldsymbol{x} = \left(3, 1\right)^\top$ 分解为行空间分量与零空间分量之和，即写出 $\boldsymbol{x} = \boldsymbol{x}_{R} + \boldsymbol{x}_{N}$。
\end{enumerate}

\begin{proof}
    对 $A$，施 Gauss - Jordan 消元法有
    \[
        A = \begin{bmatrix}
            1 & 2 \\
            2 & 4 \\
        \end{bmatrix} \implies R = \begin{bmatrix}
            1 & 2 \\
            0 & 0 \\
        \end{bmatrix}.
    \]
    这说明
    \[
    \mathrm{rank} \left(A\right) = 1,
    \]
    并且 $A$ 的列空间有一组基
    \[
        C \left(A\right) = \mathrm{span} \left\{ \left(1, 2\right)^\top \right\},
    \]
    并且 $A$ 的零空间有一组基
    \[
        N \left(A\right) = \mathrm{span} \left\{ \left(-2, 1\right)^\top \right\}.
    \]
    为求 $A$ 的 SVD，计算特征值
    \[
        \det \left(\lambda I - A^\top A\right) = \det\begin{bmatrix}
            \lambda - 5 & -10 \\
            -10 & \lambda - 20 \\
        \end{bmatrix} = 0 \implies \lambda_1 = 25, \lambda_2 = 0 \implies \sigma_1 = \sqrt{25} = 5.
    \]
    对应得到分解
    \[
        A = U \Sigma V^\top = \begin{bmatrix}
            \frac{1}{\sqrt{5}} & -\frac{2}{\sqrt{5}} \\
            \frac{2}{\sqrt{5}} & \frac{1}{\sqrt{5}} \\
        \end{bmatrix} \begin{bmatrix}
            5 & 0 \\
            0 & 0 \\
        \end{bmatrix} \begin{bmatrix}
            \frac{1}{\sqrt{5}} & -\frac{2}{\sqrt{5}} \\
            \frac{2}{\sqrt{5}} & \frac{1}{\sqrt{5}} \\
        \end{bmatrix}^\top.
    \]
    从而有 $U, V$ 与 $A$ 的四个子空间的关系
    \[
    \begin{aligned}
        N \left(A\right) &= \mathrm{span} \left\{\boldsymbol{v}_2\right\}, \\
        C \left(A^\top\right) &= \mathrm{span} \left\{\boldsymbol{v}_1\right\}, \\
        N \left(A^\top\right) &= \mathrm{span} \left\{\boldsymbol{u}_2\right\}, \\
        C \left(A\right) &= \mathrm{span} \left\{\boldsymbol{u}_1\right\}. \\
    \end{aligned}
    \]
    对 $\boldsymbol{x} = \left(3, 1\right)^\top$，行空间由右奇异向量 $\boldsymbol v_1$ 张成，故
    \[
        \boldsymbol{x}_{R} = \frac{\boldsymbol{x}^\top \boldsymbol{v}_1}{\boldsymbol{v}_1^\top \boldsymbol{v}_1} \boldsymbol{v}_1 = \frac{3 \cdot \frac{1}{\sqrt{5}} + 1 \cdot \frac{2}{\sqrt{5}}}{1} \boldsymbol{v}_1 = \sqrt{5} \boldsymbol{v}_1, \quad \boldsymbol{x}_{N} = \boldsymbol{x} - \boldsymbol{x}_{R},
    \]
    即 \[
        \boldsymbol{x}_{R} = \left(1, 2\right)^\top, \quad \boldsymbol{x}_{N} = \left(2, -1\right)^\top.
    \]
\end{proof}

\section*{Problem 2}

设矩阵
\[
A =
\begin{bmatrix}
    0 & -1 \\
    2 & 0 \\
\end{bmatrix}.
\]
\begin{enumerate}
    \item 对 $A$ 进行 SVD，求 $U, \Sigma, V$。
    \item 由 SVD 写出极分解 $A = QS$，其中 $Q = UV^\top$，$S = V \Sigma V^\top$，明确给出 $Q$ 和 $S$ 的数值。
    \item 验证 $S$ 是对称半正定矩阵，并说明 $S$ 的特征值与 $A$ 的奇异值有何关系。
\end{enumerate}

\begin{proof}
    对 $A$，计算奇异值
    \[
        \det \left(\lambda I - A^\top A\right) = \det\begin{bmatrix}
            \lambda - 4 & 0 \\
            0 & \lambda - 1 \\
        \end{bmatrix} = 0 \implies \lambda_1 = 4, \lambda_2 = 1 \implies \sigma_1 = 2, \sigma_2 = 1.
    \]
    从而有 $A$ 的 SVD
    \[
        A = U \Sigma V^\top = \begin{bmatrix}
            0 & -1 \\
            1 & 0 \\
        \end{bmatrix} \begin{bmatrix}
            2 & 0 \\
            0 & 1 \\
        \end{bmatrix} \begin{bmatrix}
            1 & 0 \\
            0 & 1 \\
        \end{bmatrix}^\top.
    \]
    从而有极分解
    \[
        A = QS = UV^\top V\Sigma V^\top = U\Sigma V^\top,
    \]
    即
    \[
        Q = UV^\top = \begin{bmatrix}
            0 & -1 \\
            1 & 0 \\
        \end{bmatrix}, \quad S = V\Sigma V^\top = \begin{bmatrix}
            2 & 0 \\
            0 & 1 \\
        \end{bmatrix}.
    \]
    验证 $S$ 是对称半正定矩阵，计算 $S$ 的特征值
    \[
        \det \left(\lambda I - S\right) = \det\begin{bmatrix}
            \lambda - 2 & 0 \\
            0 & \lambda - 1 \\
        \end{bmatrix} = 0 \implies \lambda_1 = 2, \lambda_2 = 1 > 0,
    \]
    从而 $S$ 的特征值为 $1, 2$，即 $S$ 的特征值与 $A$ 的奇异值相同。
\end{proof}

\section*{Problem 3}

设矩阵
\[
A = \begin{bmatrix}
    0 & 2 \\
    0 & 0 \\
\end{bmatrix}, \quad A_\varepsilon = \begin{bmatrix}
    0 & 2 \\
    \varepsilon & 0 \\
\end{bmatrix}, \quad \varepsilon > 0.
\]
\begin{enumerate}
    \item 计算 $A$ 的全部特征值；再计算 $A_\varepsilon$ 的特征值，说明当 $\varepsilon$ 很小时特征值的变化情况。
    \item 分别计算 $A^\top A$ 与 $A_\varepsilon^\top A_\varepsilon$ 的特征值，写出 $A$ 与 $A_\varepsilon$ 各自的奇异值，比较二者的差异。
    \item 综合以上结果，说明奇异值相比特征值具有更好数值稳定性的原因。
\end{enumerate}

\begin{proof}
    计算 $A$ 的特征值
    \[
        \det \left(\lambda I - A\right) = \det\begin{bmatrix}
            \lambda & -2 \\
            0 & \lambda \\
        \end{bmatrix} = 0 \implies \lambda_1 = 0, \lambda_2 = 0.
    \]
    计算 $A_\varepsilon$ 的特征值
    \[
        \det \left(\lambda I - A_\varepsilon\right) = \det\begin{bmatrix}
            \lambda & -2 \\
            -\varepsilon & \lambda \\
        \end{bmatrix} = 0 \implies \lambda_1 = \sqrt{2\varepsilon}, \lambda_2 = -\sqrt{2\varepsilon}.
    \]
    当 $\varepsilon$ 很小时，$A_\varepsilon$ 的特征值从 $0$ 分裂为 $\pm\sqrt{2\varepsilon}$；其变化量是 $O(\sqrt{\varepsilon})$，大于扰动本身的 $O(\varepsilon)$。

    计算 $A^\top A$ 的特征值
    \[
        A^\top A = \begin{bmatrix}
            0 & 0 \\
            2 & 0 \\
        \end{bmatrix} \begin{bmatrix}
            0 & 2 \\
            0 & 0 \\
        \end{bmatrix} = \begin{bmatrix}
            0 & 0 \\
            0 & 4 \\
        \end{bmatrix},
    \]
    从而 $A$ 的奇异值为 $\sigma_1 = 2, \sigma_2 = 0$。计算 $A_\varepsilon^\top A_\varepsilon$
    \[
        A_\varepsilon^\top A_\varepsilon =
        \begin{bmatrix}
            0 & \varepsilon \\
            2 & 0
        \end{bmatrix}
        \begin{bmatrix}
            0 & 2 \\
            \varepsilon & 0
        \end{bmatrix}
        =
        \begin{bmatrix}
            \varepsilon^2 & 0 \\
            0 & 4
        \end{bmatrix},
    \]
    因此当 $0<\varepsilon<2$ 时，$A_\varepsilon$ 的奇异值为 $\sigma_1=2,\sigma_2=\varepsilon$。与 $A$ 的奇异值 $2,0$ 相比，变化量为 $O(\varepsilon)$，与矩阵扰动同阶；这说明奇异值相较于非正规矩阵的特征值具有更好的数值稳定性。
\end{proof}
